% frontmatter/slowa/mapa-del-curso.tex -- el año entero en una tabla,
% en polaco. La escalera de sílabas de la columna "Jakie sylaby" es la
% misma que comprueba tools/gen_palabras.py (ESCALERA_PL): si cambia
% allí, cambia aquí.
\section*{Mapa roku}

\begin{center}
\renewcommand{\arraystretch}{1.4}
\begin{tabularx}{\linewidth}{@{}c c c >{\raggedright\arraybackslash}p{27mm} >{\raggedright\arraybackslash}X@{}}
\toprule
\textbf{Część} & \textbf{Dni} & \textbf{Pora roku} & \textbf{Co się czyta} & \textbf{Jakie sylaby} \\
\midrule
1 & 1--65    & jesień & 1 słowo dziennie  & Tylko sylaby otwarte, z jedną literą przed samogłoską: \emph{ma,
ko, ło, ty}. Najpierw słowa z dwiema sylabami (\emph{ma·ma, ko·ło});
od tygodnia 6 także z trzema (\emph{ło·pa·ta}). \\
2 & 66--130  & zima   & 2 słowa dziennie  & Przychodzą sylaby zamknięte (\emph{kot, lal·ka}), miękkie
(\emph{zi·ma, pies, koń}) i litery \emph{ó} i \emph{ż} (\emph{lód,
ża·ba}). \\
3 & 131--195 & wiosna & 3 słowa dziennie  & Już wszystkie: dwuznaki (\emph{sz, cz, rz, ch, dz}:
\emph{szko·ła, rze·ka}), \emph{ą} i \emph{ę} (\emph{rę·ka}) i dwie
spółgłoski przed samogłoską (\emph{kro·wa, słoń}). Słowa do czterech
sylab. \\
4 & 196--260 & lato   & 1 zdanie dziennie & Już nie ma pojedynczych słów: krótkie zdanie, od dwóch słów
(\emph{Toby czeka.}) do czterech (\emph{Dani czyta piękne bajki.}). To
most do zeszytu ze zdaniami. \\
\bottomrule
\end{tabularx}
\end{center}

\vspace{6mm}
Każda część ma 13 tygodni po 5 dni -- także w wakacje: jedna strona na
każdy dzień od poniedziałku do piątku, w Boże Narodzenie i latem tak
samo jak w roku szkolnym. Każdy tydzień opowiada kawałek historii
Lucíi, Daniego i Toby'ego, z tymi samymi tematami, tydzień po
tygodniu, co hiszpański zeszyt \emph{Aprendo a leer}: jeśli w domu jest
rodzeństwo, każde może mieć swój zeszyt i w danym tygodniu oboje czytają
o tym samym.

\vspace{4mm}
\textbf{Lucía}, \textbf{Dani} i \textbf{Toby} to słowa, które od
pierwszego dnia czyta się „od razu”, całe, jak własne imię -- chociaż
ich sylaby jeszcze nie przyszły, albo nie są polskie (hiszpańskie \emph{í}
w imieniu Lucíi; \emph{ni} w imieniu Daniego, które po polsku byłoby
miękkie; \emph{y} w imieniu Toby'ego, które brzmi tu jak \emph{i}). To imiona bohaterki, jej brata i ich psa, i pojawiają
się przez cały rok.

\vspace{4mm}
Na końcu każdej części jest medal, a na końcu zeszytu dyplom i
\textbf{odpowiedzi} (do zadań \lblAdivina, \lblUne, \lblEncuentra\ i
\lblSiNo), gdyby trzeba było coś sprawdzić.
\newpage
